\documentstyle[12pt,suthesis]{report}
\setcounter{tocdepth}{1}
\setlength{\unitlength}{1cm}

\pagestyle{empty}

\begin{document}

\begin{center}
\LARGE

${\cal YAP}$

\vspace{3cm}

A program for 3D grid visualization

\large

\vspace{7cm}

October 1993

\vspace{1 cm}

Manu Schnetzler

\vspace{2cm}

Stanford Center for Reservoir Forecasting

Stanford University, California

\end{center}

\clearpage
\pagenumbering{arabic}
\pagestyle{headings}

\tableofcontents

\listoffigures

\clearpage

\chapter{User's manual}

\section{Purpose of the program}

YAP is a program for display of two-dimensional sections ({\em slices})
of three dimensional rectangular grids. It is written in C using
the Motif library and is menu-driven.
The input file is an ascii file in the GSLIB format.
Files in binary format are also accepted to speed up the input loading (but
the BINARY format is recognized only by YAP).

In addition to on-screen visualization of the 2D sections, it
is possible to produce a postscript output of any section.

The color or gray scales ({\em colormaps}) are controlled by the user.

\section{Requirements -- compilation}
\label{compilation}

A C compiler and the Motif library are necessary to compile the program.
A {\em README} file provides the information needed to do this.

A Makefile is provided, which must be customized to fit
the architecture and location of libraries and include files.
Modify also the file {\em yaconfig.h} to define the path to access the default
colormaps, help file and some other default values
(window size, default trimming limits).

{\em make} will build YAP.

{\em make install} will build YAP and place it in the directory defined in the
Makefile ({\em EXEDIR}).

The program has been tested on ULTRIX, OSF1, SGI, HP-UX and AIX operating
systems.

\section{File formats}
\label{formats}

 \subsection{GSLIB format}
 \label{gslib_format}

A ``GSLIB-format'' file is a text ({\em ASCII}) file containing the
following information, see {\em Deutsch and Journel (1992) p.19}:
\begin{verbatim}
Title
Number of variables n
Name of variable_1
Name of variable_2
...
Name of variable_n
value_1 value_2 ... value_n
value_1 value_2 ... value_n
...
value_1 value_2 ... value_n
\end{verbatim}

The title  and names of variables are strings ending with an {\em End-Of-Line}
The number of variable is an integer.
The number of values on each line, and the number of variable names, must
match the number of variables. The values are of real (float) or integer type.
{\em They cannot be strings.}

In case of a 3-D gridded file, the GSLIB convention applies:
$x$ cycles fastest, then $y$ and finally $z$.

A GSLIB file does not include information about the grid dimensions
(number of nodes in x, y and z). Therefore, the user must provide that
information.
An ``extended'' format provides this information on the second line of the
file, after the number of variables.
Such a file is still a valid GSLIB file in the sense that it is readable
by GSLIB programs (though these programs won't use this information).

Example: for a grid $x_1, x_2,... x_k$ by $y_1, y_2,... y_m$ by
$z_1, z_2,... z_n$ with only one variable, the file would contain: \\
\\
Title \\
1 $k$ $m$ $n$ \\
Variable name \\
value at $x_1, y_1, z_1$ \\
value at $x_2, y_1, z_1$ \\
... \\
value at $x_k, y_1, z_1$ \\
value at $x_1, y_2, z_1$ \\
... \\
value at $x_k, y_m, z_1$ \\
value at $x_1, y_1, z_2$ \\
... \\
value at $x_k, y_m, z_n$ \\

The BINARY format input file contains the grid definition as well
as the minimum and maximum values of the variable. Two programs called
{\em gsl2fbin} and {\em gsl2cbin} are provided to make the conversion between
GSLIB and BINARY formats outside YAP, see next section ~\ref{gsl2bin}.

\clearpage
\subsection{Colormap files}

Colormap files follow the format:

\begin{verbatim}
Number of colors (excluding the out-of-range colors)
r g b for values < minimum
r g b for color 1
...
r g b for color n
r g b for values > maximum
\end{verbatim}

Each line gives the values of red, green and blue, between 0 and 1 (real
number). For example: \\
0 0 0 defines the color black \\
1 1 1 defines the color white \\
1 0 0 defines the color red

A grayscale from black to white, with extreme values (larger than the maximum
and lower than the minimum to display) in white could be specified as:

\begin{verbatim}
2
1 1 1    (white for values > maximum to display)
0 0 0    (black ...)
1 1 1          (... to white)
1 1 1    (white for values < minimum to display)
\end{verbatim}

\section{Conversion between formats}
\label{gsl2bin}

Though it is possible to save a GSLIB file in binary format (see
section~\ref{saving}), it might be
convenient to make the conversion outside of YAP (loading the data
file is then much faster).

Compile the programs by {\em make gsl2fbin} and {\em make gsl2cbin}.

Run the programs by
\begin{itemize}
 \item {\em gsl2fbin input\_file variable\_no nx ny nz output\_file} to create a
       BINARY float format file.
 \item {\em gsl2cbin input\_file variable\_no nx ny nz output\_file} to create a
       BINARY char format file.
\end{itemize}

The variable number specifies which variable must be saved in the binary file
(only one can be saved). This number has to be between 1 and the number
of variables in the file.

\section{Troubleshooting}

This program is an academic product, not a commercial product supported
by 800-numbers.
Please report any bug or major improvement to the author,
by email at {\em marsu@pangea.stanford.edu}, or A.G. Journel, at SCRF, by
fax: (415) 725-2099, or by email at {\em journel@pangea.stanford.edu}.

Some already known bugs and quirks:

\begin{itemize}
 \item If you close a 2D window by the window manager menu (alt-F4),
the program does not know the window has been closed and you will not be able
to open it again.
 \item If you ask for too many color classes, the allocation may fail.
In this case, the program will not crash but it will only allocate the number
of colors it is able to. If many colors are already used by other programs
(background for example), the number of remaining colors may be low.
 \item Be sure the format of your file is correct. YAP may not
recognize some errors and produce strange results or, in the best case, crash.
 \item If a 2D slice does not appear, even though the 2D window is there,
try to increase the size of this window.
\end{itemize}

To get more information about what is happening during the execution
of the program, you can open an information window in the menu {\em Edit}.
The information given there can help trace the source of the
problem.

The source code is available by anonymous ftp at {\em banach.stanford.edu}
in the directory {\em graphlib/yap}.

\clearpage

% \appendix

% \addcontentsline{toc}{section}{Appendix}

% \begin{center}
% \LARGE
% Appendix
% \vspace{2cm}
% \end{center}

\chapter{Example of run}

This section describes step by step a complete session of YAP.

\section{Starting the program}

Assuming that you have successfully compiled YAP, and that it is accessible
(either the program is in the current directory or accessible through PATH),
run the program by {\em yap}.

It is also possible to give the file name on
the command line, as well as some options. Enter {\em yap -h} for a description
of the command line arguments.

\section{Loading a file}

When YAP is started, only a white drawing area and a menu bar appear. \\
% \begin{figure}[ht]
\vspace{33pt}
\special{psfile=menubar.ps hoffset=0 voffset=-33pt}
\addcontentsline{lof}{figure}{Window -- Menu Bar}
% \end{figure}

\clearpage
Select {\em File-Open} to load a data file.

A file selection box appears to help you choose the file.\\
% \begin{figure}
\vspace{377pt}
\special{psfile=file_sel.ps hoffset=0 voffset=-377pt}
\addcontentsline{lof}{figure}{Window -- File Selection}
% \end{figure}

\clearpage

Once the file is selected (by double-clicking on the file name or clicking on the
file name and then on {\em OK}), a new window appears allowing you to choose
the variable you want to display and the trimming limits.

{\bf Remark}: in case of a binary file, this window does not appear, since
such file contains only one variable.\\
% \begin{figure}[ht]
\vspace{194pt}
\special{psfile=variables.ps hoffset=0 voffset=-194pt}
\addcontentsline{lof}{figure}{Window -- Variables Selection}
% \end{figure}

\clearpage

Click on {\em Set} and give the number of pixels in each direction, the
number of classes (not too big!), and the minimum and maximum values to
display.\\
% \begin{figure}[ht]
\vspace{346pt}
\special{psfile=parameters.ps hoffset=0 voffset=-346pt}
\addcontentsline{lof}{figure}{Window -- Parameters Selection}
% \end{figure}

The 3D volume and the color scale appear in the main window.

\section{Changing the colormap}

Select {\em File-Load colormap} to open a file selection box displaying
the file names with {\em cmap} as extension. Choose the colormap.
A default colormap (variable {\em \_CMAP\_DEF} defined in the file
{\em yaconfig.h}) is loaded if
no other colormap has been loaded before the first data file is read.
In case the default colormap is not defined or found, a gray scale is chosen.

The default location of the colormap file is defined either as an
environmental variable called {\em CMAP\_DIR} or in the file {\em yaconfig.h}
(variable {\em \_CMAP\_DIR}).
If neither of these variables are defined, the program will look in
the current directory for the colormaps.

\section{Changing the cutoffs}

You can change the cutoffs of the color scale by:
\begin{itemize}
 \item Selecting {\em Edit-Number of classes, min/max}. This gives you the
       opportunity to change the number of classes (excluding the two
       out-of-range classes) and minimum and maximum to display. The
       cutoffs are then interpolated linearly.\\
% \begin{figure}[ht]
\vspace{189pt}
\special{psfile=nb_class.ps hoffset=0 voffset=-189pt}
\addcontentsline{lof}{figure}{Window -- Number of Classes}
% \end{figure}
 \item Selecting {\em Edit-Cutoffs values}. Each cutoff value can be changed
       individually. The set of all cutoff values is accepted only if the
       cutoffs are all increasing.\\
% \begin{figure}[ht]
\vspace{160pt}
\special{psfile=cutoffs.ps hoffset=0 voffset=-160pt}
\addcontentsline{lof}{figure}{Window -- Cutoffs Values}
% \end{figure}
\end{itemize}

\section{Opening a 2D window}

You can now open a 2D window in each of the three planes (yz, xz or xy)
with the submenus in {\em 2d slices}.\\
% \begin{figure}
\vspace{358pt}
\special{psfile=2dwin.ps hoffset=0 voffset=-358pt}
\addcontentsline{lof}{figure}{Window -- 2D Window}
% \end{figure}

{\bf Remark}: the 2D views are refreshed each time the window is exposed
after having been hidden, or each time it is resized.
The 3D view, however, is refreshed only in 
case of resizing. To force a refresh, click on {\em Refresh 3D} or
{\em Refresh 2D}.

In each 2D window, you can change the slice currently shown with the slider
(except if there is only one plane in this direction).
A 2D window can be closed, refreshed individually, or sent to a postscript
file. Clicking on a 2D slice gives the information of the location (x, y and z)
and the value at this location.

\section{Printing}
When you select {\em Print} in a 2D window, you can specify:
\begin{itemize}\itemsep -0.3cm
\item name of the output postscript file
\item main title
\item titles for the two axis
\item offsets of the postscript picture (x and y positions for the beginning
      of the image, in points)
\item size of the postscript picture (in points)
\item y/x ratio (ratio of the pixels size)
\item mode (landscape or portrait)
\item if you want the color scale to appear on the image
\end{itemize}

These parameters are set in a window:\\
% \begin{figure}[ht]
\vspace{389pt}
\special{psfile=print.ps hoffset=0 voffset=-389pt}
\addcontentsline{lof}{figure}{Window -- Print}
% \end{figure}

The meaning of the parameters is given in the following drawing:\\
\vspace{423pt}
\special{psfile=xfig1.ps hoffset=0 voffset=-423pt}
\addcontentsline{lof}{figure}{Print parameters}

{\bf Remark}: the postscript output is a color image. However, if the file
is sent to a black and white postscript printer, there is an automatic
detection and the file is printed with a black to white scale.
In case the detection is not automatic with your printer, uncomment the line
\begin{verbatim}
% /color false def
\end{verbatim}
by removing the \% sign

\section{Saving a file}
\label{saving}

You can save a file by selecting {\em File-Save as}, either in GSLIB
format (with only one column) or binary (much faster to load, but recognized
only by YAP).
The BINARY char format creates smaller files. It should be used only with
integer data between 0 and 255.\\
% \begin{figure}[ht]
\vspace{184pt}
\special{psfile=save.ps hoffset=0 voffset=-212pt}
\addcontentsline{lof}{figure}{Window -- Save}
% \end{figure}

\clearpage
\section{Getting help}

If you have problems compiling the program, refer to section \ref{compilation}:
{\em Requirements -- compilation}.
Check everything is correct in the Makefile (location of libraries, includes,
compiler).

If YAP does not find the colormap files, check the directories
defined in {\em yaconfig.h}.

If you need information during the execution, the {\em Help} menu provides
limited help.\\
% \begin{figure}[ht]
\vspace{278pt}
\special{psfile=help.ps hoffset=-30pt voffset=-278pt}
\addcontentsline{lof}{figure}{Window -- Help}
% \end{figure}

\clearpage

An information windows tells you what the program does during its execution
(Menu {\em Edit-Information window}).\\
% \begin{figure}[ht]
\vspace{212pt}
\special{psfile=info_win.ps hoffset=-30pt voffset=-212pt}
\addcontentsline{lof}{figure}{Window -- Information Window}
% \end{figure}

\clearpage
\chapter{Figures}

The following pages show some screen view during the execution of the program
and postscript outputs.

\begin{figure}[h]
\vspace{560pt}
\special{psfile=screen1.ps hoffset=-220pt voffset=0pt}

\begin{center}
\vspace{20pt}
Screen view -- Continuous variable
\end{center}

\addcontentsline{lof}{figure}{Screen view -- Continuous variable}
\end{figure}

\begin{figure}[h]
\vspace{560pt}

\special{psfile=screen2.ps hoffset=-50pt voffset=0pt}

\begin{center}
\vspace{20pt}
Screen view -- Indicator variable (facies)
\end{center}

\addcontentsline{lof}{figure}{Screen view -- Indicator variable (facies)}
\end{figure}

\begin{figure}[h]
\vspace{600pt}

\begin{center}
Output of YAP -- Indicator variable (facies)
\end{center}

\special{psfile=ps_screen2.ps hoffset=-100pt voffset=0pt}
\addcontentsline{lof}{figure}{Output of YAP -- Indicator variable (facies)}
\end{figure}

\end{document}
